% ======================================================================
% Appendix K. Scattering determinant and Maass--Selberg relations
% PATH: src/appendices/K-scattering-determinant-and-maass-selberg.tex
% ======================================================================

\appendix
\chapter{Scattering determinant and Maass--Selberg relations}
\label{app:K.scattering}

\section{Purpose and interface with the main text}
\label{sec:K.purpose}

This appendix records the analytic input needed to control the continuous-spectrum sector in the localized trace
framework.  In noncompact finite-volume settings the trace formula contains an unavoidable scattering contribution,
typically expressed in terms of the logarithmic derivative of a scattering determinant.  The Maass--Selberg relations
provide the canonical mechanism by which truncation procedures are shown to be stable and by which the continuous
spectrum enters the trace identity with explicit renormalization terms.

Our goal is not to reprove the full scattering theory, but to isolate the exact statements required for the error
ledger introduced in Chapter~\ref{chap:08-normalization}.  In particular, we identify:

\begin{enumerate}
\item the normalized scattering determinant $\varphi(s)$ and its analytic continuation;
\item the Maass--Selberg renormalization identity controlling truncated Eisenstein inner products;
\item the induced trace contributions and the canonical subtraction terms;
\item the corresponding remainder bounds that enter the consolidated ledger
$\mathcal{E}_\Gamma^{\mathrm{cont}}[\cdot]$.
\end{enumerate}

Throughout we adopt the general spectral setup fixed in Chapters~\ref{chap:02-preliminaries}--%
\ref{chap:07-main-results}.  In particular, $X=\Gamma\backslash\mathbb{H}$ denotes a finite-area hyperbolic surface
(with cusps), and $\Delta$ is the nonnegative Laplace--Beltrami operator on $L^2(X)$.

\section{Eisenstein series and the scattering matrix}
\label{sec:K.eisenstein}

\subsection{Cusps and normalization}
\label{subsec:K.cusps}

Let $X=\Gamma\backslash\mathbb{H}$ be a finite-area hyperbolic surface with $\kappa\ge 1$ cusps.
Fix a complete set of cusp representatives $\{\mathfrak{a}_1,\dots,\mathfrak{a}_\kappa\}$.
For each cusp $\mathfrak{a}$ choose a scaling matrix $\sigma_{\mathfrak{a}}\in \PSL(2,\R)$ such that
$\sigma_{\mathfrak{a}}\infty=\mathfrak{a}$ and
\[
\sigma_{\mathfrak{a}}^{-1}\Gamma_{\mathfrak{a}}\sigma_{\mathfrak{a}}
=
\left\{
\pm
\begin{pmatrix} 1 & n \\ 0 & 1 \end{pmatrix}
: n\in\Z
\right\}.
\]
In the cusp coordinate $z_{\mathfrak{a}}=\sigma_{\mathfrak{a}}^{-1}z=x_{\mathfrak{a}}+iy_{\mathfrak{a}}$,
the height function is $y_{\mathfrak{a}}(z)=\Im(\sigma_{\mathfrak{a}}^{-1}z)$.

\subsection{Eisenstein series}
\label{subsec:K.eisenstein.def}

For $\Re(s)>1$ define the Eisenstein series associated with cusp $\mathfrak{a}$ by
\begin{equation}\label{eq:K.eisenstein.def}
E_{\mathfrak{a}}(z,s)
:=
\sum_{\gamma\in \Gamma_{\mathfrak{a}}\backslash\Gamma}
\Im(\sigma_{\mathfrak{a}}^{-1}\gamma z)^s.
\end{equation}
It is classical that $E_{\mathfrak{a}}(z,s)$ converges absolutely for $\Re(s)>1$,
extends meromorphically to $s\in\C$, and satisfies the eigenvalue equation
\begin{equation}\label{eq:K.eisenstein.eig}
\Delta E_{\mathfrak{a}}(\cdot,s)
=
s(1-s)\,E_{\mathfrak{a}}(\cdot,s).
\end{equation}

\subsection{Constant term and scattering matrix}
\label{subsec:K.constant.term}

Fix cusps $\mathfrak{a},\mathfrak{b}$.
The Fourier expansion of $E_{\mathfrak{a}}(z,s)$ in the cusp coordinate at $\mathfrak{b}$ has constant term
\begin{equation}\label{eq:K.constant.term}
E_{\mathfrak{a}}(\sigma_{\mathfrak{b}} z,s)
=
\delta_{\mathfrak{a}\mathfrak{b}}\,y^s
+
\phi_{\mathfrak{a}\mathfrak{b}}(s)\,y^{1-s}
+
\text{(nonconstant Fourier modes)}.
\end{equation}
The matrix
\[
\Phi(s)=(\phi_{\mathfrak{a}\mathfrak{b}}(s))_{1\le \mathfrak{a},\mathfrak{b}\le \kappa}
\]
is called the \emph{scattering matrix}.  It is meromorphic on $\C$ and satisfies the functional identity
\begin{equation}\label{eq:K.scattering.functional}
\Phi(s)\Phi(1-s)=I_\kappa.
\end{equation}

\section{Scattering determinant}
\label{sec:K.det}

\subsection{Definition and functional equation}
\label{subsec:K.det.def}

The scattering determinant is defined by
\begin{equation}\label{eq:K.det.def}
\varphi(s):=\det \Phi(s).
\end{equation}
By \eqref{eq:K.scattering.functional} it satisfies
\begin{equation}\label{eq:K.det.functional}
\varphi(s)\varphi(1-s)=1.
\end{equation}

\begin{remark}[Spectral interpretation]
\label{rem:K.det.spectral}
Poles of $\Phi(s)$ (equivalently poles of $\varphi(s)$) correspond to residual spectrum.
Zeros of $\varphi(s)$ correspond to poles of $\Phi(1-s)$ and thus encode the same residual data.
In particular, the continuous spectrum contribution to the trace formula is naturally expressed through
$\varphi'(s)/\varphi(s)$.
\end{remark}

\subsection{Logarithmic derivative bounds}
\label{subsec:K.logder.bounds}

The trace formula requires control of integrals involving $\varphi'/\varphi$ along the critical line $\Re(s)=1/2$.

\begin{assumption}[Admissible scattering growth]
\label{ass:K.scattering.growth}
There exists $C>0$ such that for $|t|\ge 2$,
\begin{equation}\label{eq:K.logder.growth}
\left|\frac{\varphi'}{\varphi}\!\left(\frac12+it\right)\right|
\le
C\log(2+|t|).
\end{equation}
\end{assumption}

\begin{remark}
\label{rem:K.scattering.growth}
Bounds of the form \eqref{eq:K.logder.growth} are standard consequences of the meromorphic continuation,
Hadamard factorization, and Weyl-type bounds for poles, and are sufficient for all smoothing arguments in the
localized trace regime.
\end{remark}

\section{Truncation and Maass--Selberg relations}
\label{sec:K.ms}

\subsection{Truncation operator}
\label{subsec:K.truncation}

Fix truncation heights $Y=(Y_1,\dots,Y_\kappa)$ with each $Y_j\ge 1$.
Let $X(Y)\subset X$ denote the truncated surface obtained by removing the cusp regions
$\{y_{\mathfrak{a}_j}>Y_j\}$.
Define the truncation operator $\Lambda^Y$ on smooth functions $f$ by
\begin{equation}\label{eq:K.truncation.op}
(\Lambda^Y f)(z)
=
f(z)
-
\sum_{j=1}^\kappa
\mathbf{1}_{y_{\mathfrak{a}_j}(z)>Y_j}
\left(
f_{\mathfrak{a}_j}^{(0)}(y_{\mathfrak{a}_j}(z))
\right),
\end{equation}
where $f_{\mathfrak{a}_j}^{(0)}(y)$ denotes the constant term of the Fourier expansion of $f$ at cusp $\mathfrak{a}_j$.

\begin{remark}
\label{rem:K.truncation}
The truncation $\Lambda^Y$ removes the divergent constant terms in the cusps, producing an $L^2$ function on $X$.
This is the canonical renormalization procedure underlying the continuous spectrum sector.
\end{remark}

\subsection{Maass--Selberg relation: matrix form}
\label{subsec:K.ms.matrix}

Let $E(z,s)$ denote the $\kappa$-vector of Eisenstein series
\[
E(z,s)=(E_{\mathfrak{a}_1}(z,s),\dots,E_{\mathfrak{a}_\kappa}(z,s))^T.
\]
For $s,w\in\C$ with $\Re(s),\Re(w)>1$ consider the truncated inner product
\[
\left\langle \Lambda^Y E(\cdot,s),\Lambda^Y E(\cdot,\overline{w})\right\rangle_{L^2(X)}.
\]

\begin{theorem}[Maass--Selberg relation]
\label{thm:K.maass.selberg}
For $s,w\in\C$ with $s\neq w$ and $s\neq 1-w$ one has the identity
\begin{align}
\label{eq:K.ms.relation}
\left\langle \Lambda^Y E(\cdot,s),\Lambda^Y E(\cdot,\overline{w})\right\rangle
&=
\frac{Y^{s+\overline{w}-1}-Y^{1-s-\overline{w}}}{s+\overline{w}-1}\,I_\kappa
\\
\nonumber
&\qquad
+
\frac{Y^{s-\overline{w}}-Y^{\overline{w}-s}}{s-\overline{w}}\,\Phi(s)
+
\text{(holomorphic remainder)},
\end{align}
where $Y^\alpha$ denotes the diagonal matrix $\diag(Y_1^\alpha,\dots,Y_\kappa^\alpha)$.
\end{theorem}

\begin{remark}[Diagonal specialization]
\label{rem:K.ms.diagonal}
Specializing $w=s$ yields the classical diagonal Maass--Selberg relation, which contains logarithmic terms in $Y$.
This diagonal identity is the one that enters directly into trace computations.
\end{remark}

\subsection{Diagonal Maass--Selberg identity}
\label{subsec:K.ms.diagonal}

We state the diagonal identity in the form required for the continuous spectrum trace term.

\begin{theorem}[Diagonal Maass--Selberg relation]
\label{thm:K.ms.diagonal}
Let $s=\sigma+it$ with $\sigma\in\R$.
Then one has
\begin{equation}\label{eq:K.ms.diagonal}
\left\langle \Lambda^Y E(\cdot,s),\Lambda^Y E(\cdot,\overline{s})\right\rangle
=
2\log Y
-
\frac{\varphi'}{\varphi}(s)
+
O(1),
\end{equation}
where $\log Y$ denotes the diagonal matrix $\diag(\log Y_1,\dots,\log Y_\kappa)$ and the $O(1)$ term is uniform in $t$
for $|t|\ge 2$ and in $Y_j\ge 1$.
\end{theorem}

\begin{remark}[Canonical content]
\label{rem:K.ms.canonical}
The term $2\log Y$ is the \emph{explicit truncation divergence}.
The logarithmic derivative $\varphi'/\varphi$ is the \emph{intrinsic scattering correction}.
The $O(1)$ term is a uniform bounded remainder that may be absorbed into the continuous-spectrum ledger.
\end{remark}

\section{Continuous spectrum contribution to the trace}
\label{sec:K.trace.cont}

\subsection{Spectral decomposition}
\label{subsec:K.spectral.decomp}

Let $h$ be an admissible test function as in Chapter~\ref{chap:08-normalization}.
In the noncompact setting, the spectral side of the trace formula contains an integral contribution
\begin{equation}\label{eq:K.cont.spectral}
\mathcal{C}[h]
=
\frac{1}{4\pi}
\int_{-\infty}^{\infty}
h(t)\,
\Tr\!\left(
-\frac{\Phi'}{\Phi}\!\left(\frac12+it\right)
\right)\,dt,
\end{equation}
where $\Tr(\Phi'/\Phi)$ denotes the trace of the matrix logarithmic derivative.

Equivalently, in determinant form,
\begin{equation}\label{eq:K.cont.det}
\mathcal{C}[h]
=
\frac{1}{4\pi}
\int_{-\infty}^{\infty}
h(t)\,
\left(
-\frac{\varphi'}{\varphi}\!\left(\frac12+it\right)
\right)\,dt.
\end{equation}

\begin{remark}[Sign convention]
\label{rem:K.sign}
The sign in \eqref{eq:K.cont.det} depends on the normalization of the Eisenstein series and on the convention for the
spectral parameter.  In the main text we fix the sign so that the continuous-spectrum term appears as a correction
added to the identity sector.  Any change of convention is absorbed into the explicit identity correction described
in Chapter~\ref{chap:08-normalization}.
\end{remark}

\subsection{Normalization compatibility}
\label{subsec:K.norm.compat}

The normalization procedure of Chapter~\ref{chap:08-normalization} removes the unstable identity-sector contribution
at $t=0$ by replacing $h$ with $h^\sharp=h-h(0)\chi$.
The continuous-spectrum term is therefore canonically defined by
\begin{equation}\label{eq:K.cont.normalized}
\mathcal{C}_\sharp[h]
:=
\frac{1}{4\pi}
\int_{-\infty}^{\infty}
h^\sharp(t)\,
\left(
-\frac{\varphi'}{\varphi}\!\left(\frac12+it\right)
\right)\,dt.
\end{equation}

\begin{lemma}[Cutoff stability of the continuous term]
\label{lem:K.cutoff.stability}
Let $\chi_1,\chi_2$ be admissible cutoffs and define $h^{\sharp,j}=h-h(0)\chi_j$.
Then
\begin{equation}\label{eq:K.cutoff.diff}
\mathcal{C}_{\sharp}^{(1)}[h]-\mathcal{C}_{\sharp}^{(2)}[h]
=
-\frac{h(0)}{4\pi}
\int_{-\infty}^{\infty}
(\chi_1(t)-\chi_2(t))
\left(
-\frac{\varphi'}{\varphi}\!\left(\frac12+it\right)
\right)\,dt,
\end{equation}
and the right-hand side is finite and explicitly computable.
\end{lemma}

\begin{proof}
Immediate from linearity and from the definition \eqref{eq:K.cont.normalized}.
\end{proof}

\section{Ledger bounds for scattering contributions}
\label{sec:K.ledger}

We now extract the quantitative statement needed for the consolidated error ledger.

\subsection{Absolute convergence under admissible tests}
\label{subsec:K.abs.conv}

\begin{lemma}[Absolute convergence]
\label{lem:K.abs.conv}
Let $h$ be admissible and assume $h$ is rapidly decaying.
Under Assumption~\ref{ass:K.scattering.growth} the integral \eqref{eq:K.cont.det} converges absolutely and
\begin{equation}\label{eq:K.abs.conv.bound}
|\mathcal{C}[h]|
\le
C
\int_{\R} |h(t)|\log(2+|t|)\,dt.
\end{equation}
\end{lemma}

\begin{proof}
Immediate from \eqref{eq:K.logder.growth}.
\end{proof}

\subsection{Continuous-spectrum ledger term}
\label{subsec:K.cont.ledger}

\begin{definition}[Continuous-spectrum ledger contribution]
\label{def:K.cont.ledger}
For admissible $h$ define
\begin{equation}\label{eq:K.cont.ledger.def}
\mathcal{E}^{\mathrm{cont}}_\Gamma[h]
:=
C_\Gamma
\int_{\R}
|h^\sharp(t)|\,\log(2+|t|)\,dt,
\end{equation}
where $C_\Gamma$ is a fixed constant dominating the implied constant in
Assumption~\ref{ass:K.scattering.growth}.
\end{definition}

\begin{remark}
\label{rem:K.cont.ledger}
The definition \eqref{eq:K.cont.ledger.def} is canonical at the level of the present thesis:
it isolates the exact integrability mechanism required for the continuous spectrum and packages the dependence into a
single stable quantity.
Sharper versions may be substituted in specific arithmetic settings where explicit bounds on $\varphi'/\varphi$ are
available.
\end{remark}

\subsection{Smoothed oscillatory tests}
\label{subsec:K.smoothed}

For the standard oscillatory family $h_{\tau,\Delta}(t)=W(t/\Delta)e^{-i\tau t}$ used in
Chapter~\ref{chap:10-trace-to-zeta}, one obtains a uniform bound of the form
\[
\mathcal{E}^{\mathrm{cont}}_\Gamma[h_{\tau,\Delta}]
\lesssim
\Delta \log(2+|\tau|),
\]
since $h_{\tau,\Delta}$ is supported in $|t|\lesssim \Delta$ and $|h_{\tau,\Delta}(t)|\le 1$.

\begin{lemma}[Continuous-spectrum bound for oscillatory windows]
\label{lem:K.osc.bound}
Let $W\in C_c^\infty(\R)$ and let $h_{\tau,\Delta}$ be defined by \eqref{eq:10.htaudelta}.
Then
\begin{equation}\label{eq:K.osc.bound}
\mathcal{E}^{\mathrm{cont}}_\Gamma[h_{\tau,\Delta}]
\le
C_{\Gamma,W}\,\Delta\,\log(2+|\tau|).
\end{equation}
\end{lemma}

\begin{proof}
By support control and the growth bound \eqref{eq:K.logder.growth}.
\end{proof}

\section{A canonical scattering renormalization identity}
\label{sec:K.canonical.identity}

We now state the final interface lemma used implicitly in the main text.

\begin{proposition}[Canonical scattering renormalization]
\label{prop:K.canonical.renorm}
Let $h$ be admissible and let $\mathcal{T}_\Gamma[h]$ be the normalized trace functional of
Chapter~\ref{chap:08-normalization}.  In the finite-volume noncompact setting, the continuous-spectrum contribution
admits a canonical decomposition
\begin{equation}\label{eq:K.canonical.decomp}
\mathcal{C}[h]
=
\mathcal{C}_\sharp[h]
+
h(0)\,\mathcal{C}[\chi]
+
\Err_{\mathrm{cont}}[h],
\end{equation}
where:
\begin{enumerate}
\item $\mathcal{C}_\sharp[h]$ is given by \eqref{eq:K.cont.normalized};
\item the term $h(0)\,\mathcal{C}[\chi]$ is an explicit identity-sector correction depending only on the fixed cutoff;
\item the remainder satisfies
\begin{equation}\label{eq:K.cont.remainder.bound}
|\Err_{\mathrm{cont}}[h]|
\le
\mathcal{E}^{\mathrm{cont}}_\Gamma[h].
\end{equation}
\end{enumerate}
\end{proposition}

\begin{proof}
The identity follows by inserting $h=h^\sharp+h(0)\chi$ into \eqref{eq:K.cont.det} and applying the absolute bound
Lemma~\ref{lem:K.abs.conv}.
\end{proof}

\section{Consequences for the normalized trace framework}
\label{sec:K.consequences}

We record the two statements used implicitly in the main stability theorem.

\begin{corollary}[Continuous spectrum is ledger-controlled]
\label{cor:K.ledger.control}
In the noncompact finite-volume setting, all continuous-spectrum contributions to the localized trace formula are
canonical modulo explicit cutoff-dependent identity corrections, and the noncanonical remainder is controlled by
$\mathcal{E}^{\mathrm{cont}}_\Gamma[\cdot]$.
\end{corollary}

\begin{corollary}[Compatibility with the consolidated ledger]
\label{cor:K.consolidated}
With $\mathcal{E}^{\mathrm{cont}}_\Gamma[\cdot]$ defined as in \eqref{eq:K.cont.ledger.def}, the total consolidated
ledger of Chapter~\ref{chap:08-normalization} remains valid without modification in the presence of a continuous
spectrum.
\end{corollary}

\section*{Bibliographic notes}

The scattering matrix, scattering determinant, and Maass--Selberg relations are classical components of the spectral
theory of finite-area hyperbolic surfaces.  Standard references include the Selberg trace formula treatments of
Hejhal \cite{Hejhal1983} and the spectral methods exposition of Iwaniec \cite{Iwaniec2002}.  For systematic analytic
scattering theory in geometric contexts, see also \cite{DyatlovZworski2019}.

% ======================================================================
% End of K-scattering-determinant-and-maass-selberg.tex
% ======================================================================
